% from ega4-10-fr.tex
\oldpage[IV-3]{114}

\begin{env}[10.9.5]
\label{IV.10.9.5-fr}
Considérons maintenant un espace annelé $(X,\mathcal O_X)$~; nous dirons qu'une
partie ouverte $U$ de $X$ est \emph{ultra-affine} si l'espace annelé induit
$(U,\mathcal O_X|U)$ est isomorphe à un spectre maximal
$\operatorname{Spm}(A)$, où $A$ est un anneau de Jacobson. Nous dirons que $X$
est un \emph{ultra-préschéma} si tout point de $X$ admet un voisinage ouvert
ultra-affine. On montre, comme dans (\textbf I, 2.1.3 et 2.1.4) que les ouverts
ultra-affines forment une \emph{base} de la topologie de $X$ et que $X$ est un
espace de Kolmogoroff. Si $Y$ est un second ultra-préschéma, nous dirons qu'un
morphisme d'espaces annelés $f:X\to Y$ est un \emph{morphisme
d'ultra-préschémas} s'il vérifie les conditions suivantes~: $1^\circ$ $f$ est
un morphisme d'espaces annelés en anneaux locaux~; $2^\circ$ pour tout $x\in X$,
il y a un voisinage ouvert ultra-affine $V$ de $f(x)$ dans $Y$ et un voisinage
ouvert ultra-affine $U$ de $x$ dans $X$ tels que $f(U)\subset V$ et que
l'homomorphisme
$\Gamma(V,\mathcal O_Y)\to\Gamma(U,\mathcal O_X)$ correspondant à $f$ fasse
de $\Gamma(U,\mathcal O_X)$ une $\Gamma(V,\mathcal O_Y)$-algèbre de type fini.

Il est immédiat qu'on définit bien ainsi des morphismes, le composé de deux
morphismes en étant un troisième grâce à la remarque finale de (10.9.3). Il
est clair que la catégorie $\mathbf C''_0$ ainsi définie est une sous-catégorie
de $\mathbf C'$ qui contient $\mathbf C''$~; nous nous proposons de montrer que
$\mathbf C''=\mathbf C''_0$, autrement dit~:
\end{env}

\begin{proposition}[10.9.6]
\label{IV.10.9.6-fr}
Le foncteur $X\mapsto S(X)$ de $\mathbf C$ dans $\mathbf C''_0$ est une
équivalence de catégories.
\end{proposition}

$1^\circ$ Montrons d'abord que le foncteur $X\mapsto S(X)$ est
\emph{pleinement fidèle}, autrement dit que pour $X$, $Y$ dans $\mathbf C$,
l'application canonique
\[
  \operatorname{Hom}_{\mathbf C}(X,Y)\longrightarrow
  \operatorname{Hom}_{\mathbf C''_0}(S(X),S(Y))
\]
est \emph{bijective}. En premier lieu elle est \emph{injective}~: soient en
effet $f$, $g$ deux morphismes localement de type fini de $X$ dans $Y$ et
supposons que $S(f)=S(g)$. Cela entraîne d'abord que pour tout ouvert $V$ de
$Y$, on a
$f^{-1}(V)\cap S(X)=g^{-1}(V)\cap S(X)$, donc (10.3.1 et 10.2.7)
$f^{-1}(V)=g^{-1}(V)$~; il suffit donc de prouver que pour tout ouvert affine
$V$ de $Y$, $f$ et $g$ coïncident dans
$f^{-1}(V)=g^{-1}(V)$, autrement dit, on est ramené au cas où
$Y=\operatorname{Spec}(B)$ est le spectre d'un anneau de Jacobson $B$. Il
suffit (\textbf I, 2.2.4) de montrer que les homomorphismes d'anneaux
$B\to\Gamma(X,\mathcal O_X)$ correspondant à $f$ et $g$ sont alors les mêmes.
Or, pour tout $s\in B$, les images de $s$ par ces homomorphismes sont deux
sections de $\mathcal O_X$ au-dessus de $X$ qui, par hypothèse, induisent la
même section au-dessus de $S(X)$~; on sait donc (10.2.8) que ces sections sont
identiques, d'où notre assertion. Prouvons en second lieu que tout morphisme
$h:S(X)\to S(Y)$ (pour la catégorie $\mathbf C''_0$) est de la forme $S(f)$,
où $f:X\to Y$ est un morphisme localement de type fini. Par hypothèse, il
existe un recouvrement ouvert ultra-affine $(U'_\lambda)$ (resp.
$(V'_\mu)$) de $S(X)$ (resp. $S(Y)$) tel que pour tout $\lambda$,
$h(U'_\lambda)$ soit contenu dans un $V'_\mu$ et qu'il corresponde au
morphisme $h_\lambda:U'_\lambda\to V'_\mu$, restriction de $h$, un
homomorphisme d'anneaux
$\Gamma(V'_\mu,\mathcal O_{S(Y)})\to
\Gamma(U'_\lambda,\mathcal O_{S(X)})$ faisant du second anneau une algèbre de
type fini sur le premier. On peut supposer que
$U'_\lambda=U_\lambda\cap S(X)$ et $V'_\mu=V_\mu\cap S(Y)$, où
$U_\lambda$ et $V_\mu$ sont des ouverts affines uniquement déterminés dans
$X$ et $Y$ respectivement~; si l'on prouve que, pour tout $\lambda$,
$h_\lambda=S(f_\lambda)$ où $f_\lambda:U_\lambda\to V_\mu$ est un morphisme
de type fini, alors il résulte de la première partie du raisonnement,
appliquée aux restrictions de $f_\alpha$ et $f_\beta$ à
$U_\alpha\cap U_\beta$, que les $f_\lambda$ sont les restrictions d'un même
morphisme $f:X\to Y$, et on aura évidemment $h=S(f)$. On est ainsi ramené au
cas où $X$ et $Y$ sont affines, et la conclusion résulte alors de (10.9.4).

$2^\circ$ Il reste à prouver que tout ultra-préschéma $X'$ est de la forme
$S(X)$ pour un préschéma de Jacobson $X$ (qui sera nécessairement unique à
isomorphisme près, en vertu de $1^\circ$). Il y a un recouvrement
$(U'_\alpha)$ de $X'$ par des ouverts ultra-affines, dont chacun est de la
forme $S(U_\alpha)$, $U_\alpha$ étant le spectre d'un anneau de Jacobson. Pour
tout couple d'indices $\alpha$, $\beta$, considérons l'unique ouvert
$U_{\alpha\beta}$ de $U_\alpha$ dont la trace sur $U'_\alpha$ est
$U'_\alpha\cap U'_\beta$~; en vertu de $1^\circ$, l'automorphisme identique de
$U'_\alpha\cap U'_\beta$ est de la forme $S(\theta_{\alpha\beta})$, où
$\theta_{\alpha\beta}:U_{\alpha\beta}\to U_{\beta\alpha}$ est un
isomorphisme de préschémas. On constate aussitôt (en vertu du $1^\circ$) que la
famille $(\theta_{\alpha\beta})$ vérifie la condition de recollement
($\mathbf 0_{\mathrm I}$, 4.1.7), et que cette famille définit donc un
préschéma $X$, tel que les $U_\alpha$ s'identifient à des ouverts affines de
$X$~; il est clair alors que l'on a $X'=S(X)$, ce qui achève la démonstration.

\subsection{Espaces algébriques de Serre}
\label{subsection:IV.10.10-fr}

\begin{env}[10.10.1]
\label{IV.10.10.1-fr}
Le langage introduit par Serre dans (FAC) est parfois commode, notamment dans
les questions où l'intérêt majeur s'attache aux points rationnels sur le corps
de base $k$ (algébriquement clos par hypothèse) des «~variétés algébriques~»
sur $k$ que l'on considère. Nous allons esquisser ici ce langage en le
rattachant aux considérations qui précèdent, pour permettre au lecteur de
traduire les énoncés de Serre en langage des schémas. Il est d'ailleurs
possible de développer le langage de Serre également pour des préschémas sur
un corps non algébriquement clos (et même sur un anneau artinien)~; mais cela
introduit des complications techniques considérables, et d'ailleurs, sur un
corps de base arbitraire, les avantages (surtout psychologiques) du point de
vue de Serre disparaissent~; aussi
